\documentclass[11pt]{article}
\usepackage[margin=1in]{geometry}
\usepackage{amsmath,amssymb}
\usepackage{booktabs}
\usepackage{hyperref}

\title{Monte Carlo Simulation of Linezolid $EC_{50}$ Variability\\
\large \texttt{MC\_ec50\_lzd.py}}
\author{}
\date{}

\begin{document}
\maketitle

\section{What it does}

\texttt{MC\_ec50\_lzd.py} runs \texttt{model\_Bacteremia.py}'s
\texttt{dual\_reservoir\_model} $N = 1000$ times from identical initial
conditions ($S_b(0){=}R_b(0){=}0$, $S_{res}(0){=}R_{res}(0){=}100$ CFU/mL) and
identical PD/PK parameters, varying only $EC_{50,L}$ -- linezolid's
potency/susceptibility parameter in the $E_{max}$ pharmacodynamic sigmoid --
once per iteration:
\begin{equation}
E_L(t) = \frac{E_{max,L}\, L(t)}{EC_{50,L} + L(t)}.
\end{equation}
$EC_{50,L}$ is sampled from a mean-corrected log-normal distribution,
\begin{equation}
\sigma_{\ln} = \sqrt{\ln(1+CV^2)}, \qquad
\mu_{\ln} = \ln(EC_{50,L,\mathrm{base}}) - \tfrac{1}{2}\sigma_{\ln}^2, \qquad
EC_{50,L} \sim \mathrm{LogNormal}(\mu_{\ln}, \sigma_{\ln}^2),
\end{equation}
with $EC_{50,L,\mathrm{base}} = 1.0$ mg/L and $CV = 1.117$, so that
$\mathbb{E}[EC_{50,L}]$ lands exactly on baseline -- the same
mean-corrected, CV-parameterized convention used by
\texttt{MC\_immune\_response.py} and \texttt{MC\_emax\_lzd.py}. Isolating
this one parameter quantifies how inter-patient variability in linezolid
susceptibility alone -- independent of dosing, immune response, or growth
fitness -- determines whether the resistant reservoir is suppressed or
escapes.

\section{Why $EC_{50,L}$ specifically, and not a fixed value}

Growth rates ($\rho_S$, $\rho_R$) were scaled $5\times$ elsewhere in the
project alongside an increase in $E_{max,L}$, but that scaling cannot
substitute for varying $EC_{50,L}$: $EC_{50,L}$ independently controls how
much linezolid effect is achieved at realistic, sub-saturating drug
concentrations (the $E_{max}$ sigmoid above), which is mechanistically
distinct from "linezolid works well for this patient" as expressed through
growth-rate or fitness-cost parameters. A patient can have a normal fitness
cost for resistance ($\rho_R/\rho_S$) yet still respond very differently to
treatment purely because their $EC_{50,L}$ places them on a different part of
the sigmoid at achievable drug levels. This script isolates that second,
independent source of variability.

\section{Fixed baseline parameters}

\begin{table}[h]
\centering
\begin{tabular}{@{}llp{7.2cm}@{}}
\toprule
Parameter & Value & Note \\
\midrule
$\rho_S$ & $0.61$ h$^{-1}$ & \\
$\rho_R$ & $0.55$ h$^{-1}$ & $\sim10\%$ fitness cost relative to $\rho_S$ \\
$\rho_{res,S}$ & $0.175$ h$^{-1}$ & \\
$\rho_{res,R}$ & $0.1765$ h$^{-1}$ & Held just above the reservoir's own persistence threshold ($0.17594055$) but below the point where $R_b$ wins the blood race outright, so $S_b$ can also establish a visible infection (peak $\sim4.8\times10^3$ CFU/mL) \\
$E_{max,V}$ & $0.40$ & \\
$EC_{50,V}$ & $0.245$ & \\
$E_{max,L}$ & $0.8$ & Fixed, decoupled from $\rho_S$ \\
$EC_{50,L,\mathrm{base}}$ & $1.0$ mg/L & Varied per MC iteration (this script's sampled parameter) \\
$B_{max,blood}$ & $6000$ CFU/mL & \\
$B_{max,res}$ & $10^4$ CFU/mL & \\
$\mathtt{van\_res\_fraction}$ & $0.15$ & \\
$\mathtt{lzd\_res\_fraction}$ & $0.45$ & \\
$f_{r\to b}$ & $5\times10^{-5}$ h$^{-1}$ & \\
$f_{b\to r}$ & $1\times10^{-5}$ h$^{-1}$ & \\
\bottomrule
\end{tabular}
\caption{\texttt{BASE\_PARAMS} in \texttt{MC\_ec50\_lzd.py}. Every value is held fixed across all 1000 iterations; only $EC_{50,L}$ varies.}
\end{table}

\section{Step-by-step procedure, as implemented}

This section follows the code's actual execution order.

\begin{enumerate}
\item \textbf{Load the model.} \texttt{model\_Bacteremia.py} is loaded
dynamically via \texttt{importlib} (rather than a plain \texttt{import}) so
a missing file produces an explicit error message instead of an
\texttt{ImportError} traceback.

\item \textbf{Fix run-level constants.} \texttt{NUM\_ITERATIONS = 1000},
\texttt{LOD = 10.0} CFU/mL, \texttt{CV = 1.117}, \texttt{SEED = 44},
\texttt{total\_h = 1944} h (21 pre-treatment days + 4-day vancomycin course
starting at \texttt{vanco\_start = 504} h + 42-day linezolid course + 14-day
post-treatment follow-up), and a shared evaluation grid
\texttt{t\_eval = linspace(0, 1944, 1150)} used by every ODE solve.

\item \textbf{Build the drug-concentration and immune-response callables
once.} \texttt{PharmacokineticModel().concentration\_function(...)} is
called once each for vancomycin and linezolid to produce \texttt{van\_func(t)}
/ \texttt{lzd\_func(t)}, and \texttt{ImmuneResponse()} is instantiated once.
These closures are reused, unmodified, by every one of the 1000 ODE solves
below -- only $EC_{50,L}$ changes between iterations.

\item \textbf{Fix \texttt{BASE\_PARAMS} and the initial state.} All
parameters in Table 1 are set once. The initial state
\texttt{Y0 = [0.0, 0.0, 100.0, 100.0]} ($S_b, R_b, S_{res}, R_{res}$) is
likewise identical for every iteration; $EC_{50,L}$ is deliberately absent
from \texttt{BASE\_PARAMS} and is instead merged in per-call (step 8).

\item \textbf{Compute the log-normal sampling parameters.}
$\sigma_{\ln} = \sqrt{\ln(1+CV^2)}$ and
$\mu_{\ln} = \ln(EC_{50,L,\mathrm{base}}) - \sigma_{\ln}^2/2$ are computed
once from $CV = 1.117$ and $EC_{50,L,\mathrm{base}} = 1.0$.

\item \textbf{Draw all 1000 samples at once.} A seeded generator,
\texttt{rng = np.random.default\_rng(44)}, draws the full vector
\texttt{ec50\_l\_samples = rng.lognormal(} $\mu_{\ln}, \sigma_{\ln},$
\texttt{1000)} in a single call -- so the run is fully reproducible given the
fixed seed, and sample $i$ is used by MC iteration $i$ in step 8.

\item \textbf{Compute the deterministic switch curve (121 points).} Before
the Monte Carlo loop, \texttt{sweep\_ec50 = linspace(0.3, 2.0, 121)} is
built, and for each swept value \texttt{\_final\_counts(ec50\_l)} solves the
ODE (\texttt{odeint}, \texttt{rtol=1e-7}, \texttt{atol=1e-9},
\texttt{mxstep=5000}) with $EC_{50,L}$ substituted into a copy of
\texttt{BASE\_PARAMS}, recording only the final-time $R_b$ and $R_{res}$
(\texttt{sol[-1, 1]}, \texttt{sol[-1, 3]}) into \texttt{sweep\_R\_b} /
\texttt{sweep\_R\_res}.

\item \textbf{Root-find \texttt{ESCAPE\_THRESH}.} On the same
\texttt{\_final\_counts} function, \texttt{scipy.optimize.brentq} solves
$R_{res}^{\mathrm{final}}(EC_{50,L}) - \mathrm{LOD} = 0$ over
$[0.3, 2.0]$, provided the sign of $f$ differs at the two endpoints; if not,
it falls back to the coarse sweep index where \texttt{sweep\_R\_res} first
exceeds the LOD.

\item \textbf{Main Monte Carlo loop (1000 iterations).} Pre-allocated
\texttt{NaN}-filled arrays (\texttt{S\_b\_hist}, \texttt{R\_b\_hist},
\texttt{S\_res\_hist}, \texttt{R\_res\_hist}, each shape
$1000 \times 1150$) are filled by, for each $i = 0, \dots, 999$:
\begin{enumerate}
\item[(a)] build \texttt{params\_i = \{**BASE\_PARAMS, "EC50\_L":
ec50\_l\_samples[i]\}} -- a fresh shallow copy of the fixed parameters with
only $EC_{50,L}$ replaced by this iteration's sample;
\item[(b)] integrate the ODE from the same fixed \texttt{Y0} over the same
\texttt{t\_eval}, with \texttt{odeint} (\texttt{rtol=1e-7},
\texttt{atol=1e-9}, \texttt{mxstep=5000});
\item[(c)] floor every one of the four state trajectories to $0$ wherever
it falls below \texttt{LOD} (\texttt{np.where(sol[:, k] < LOD, 0.0,
sol[:, k])}), so sub-detection noise reads as true clearance in later
plots/statistics rather than as a small positive count;
\item[(d)] store the full length-1150 trajectory for each compartment into
row $i$ of the four history arrays;
\item[(e)] print progress every 50 completed iterations.
\end{enumerate}
No parameter other than $EC_{50,L}$ differs between iterations -- the
100\% of the run-to-run variance in outcomes is attributable to sampled
$EC_{50,L}$ alone.

\item \textbf{Derive per-iteration summary metrics.} After the loop,
\texttt{peak\_S\_b}, \texttt{peak\_R\_b}, \texttt{peak\_S\_res} are the
per-row maxima (\texttt{np.nanmax(..., axis=1)}) of the corresponding
history arrays, and \texttt{final\_R\_b}, \texttt{final\_R\_res} are the
last time-point of each row (\texttt{hist[:, -1]}, with \texttt{NaN}
mapped to $0$).

\item \textbf{Render the three figures} (kinetics, outcomes, switch --
described in Section~\ref{sec:outputs}) directly from the history arrays,
the summary metrics, the deterministic sweep arrays, and
\texttt{ESCAPE\_THRESH}; no further simulation is performed.

\item \textbf{Print a console summary} of the sampled $EC_{50,L}$
distribution (median, 5th/95th percentile) and, where any samples exceed
the LOD, the median and 5th/95th percentile of peak $S_b$ and peak $R_b$
across iterations.
\end{enumerate}

\section{Design: deterministic switch curve, then Monte Carlo}

The script does two distinct things, in this order:

\begin{enumerate}
\item \textbf{Deterministic switch curve.} A single sweep of $EC_{50,L}$
over $[0.3, 2.0]$ mg/L (121 evenly spaced points, \texttt{SWEEP\_MIN}/
\texttt{SWEEP\_MAX}) is run first, recording final $R_b$ and final $R_{res}$
at each point. The exact $EC_{50,L}$ at which final $R_{res}$ crosses the
limit of detection (LOD $= 10$ CFU/mL) is located by root-finding
(\texttt{scipy.optimize.brentq}) on
\begin{equation}
f(EC_{50,L}) = R_{res}^{\mathrm{final}}(EC_{50,L}) - \mathrm{LOD} = 0,
\end{equation}
giving a precise escape threshold, \texttt{ESCAPE\_THRESH}. This replaces an
old hardcoded value (\texttt{ESCAPE\_THRESH} $=0.988$) that no longer
matched this file's own $\rho_S = 0.61$ baseline; the same quantity is
re-derived independently in a separate diagnostic script,
\texttt{EC50\_lzd\_threshold\_sweep.py}, which bisects the same
\texttt{BASE\_PARAMS} and finds $EC_{50,L} \approx 0.93$. This is the actual
dose--response result: a continuous curve and a single crossing point, not a
handful of coarse categorical bins.

\item \textbf{Monte Carlo sampling.} The 1000-iteration loop then asks a
different question: given realistic inter-patient $EC_{50,L}$ variability
(the log-normal distribution above), what fraction of virtual patients land
on each side of that deterministic switch, and what do their full bacterial
kinetics look like over the treatment course. Each MC sample's outcome is
colored, in the resulting figures, by its own simulated result (whether that
patient's final $R_{res}$ is above or below the LOD), not merely by which
side of \texttt{ESCAPE\_THRESH} its sampled $EC_{50,L}$ happens to fall on --
so the outcome figures stay faithful to the simulated data rather than to
the idealized deterministic switch.
\end{enumerate}

Under the current baseline, \texttt{ESCAPE\_THRESH} sits right next to
baseline ($EC_{50,L,\mathrm{base}} = 1.0$): only $\sim47\%$ of sampled
$EC_{50,L}$ values suppress the reservoir, versus $\sim53\%$ that let it
persist and relapse after treatment ends. Both outcomes are accepted as
realistic -- the purpose of the chosen $\rho_{res,R}$ is to let $S_b$
establish a visible blood infection, not to guarantee reservoir persistence
(or clearance) in every sample.

\section{Outputs}
\label{sec:outputs}

The script generates three figures:
\begin{enumerate}
\item \textbf{\texttt{mc\_ec50\_lzd\_sweep\_kinetics.png}} -- a four-panel
time-course figure ($S_b$, $R_b$, $S_{res}$, $R_{res}$) over the $1944$-h
simulation window, each panel showing 20 individual sample trajectories plus
the median and $5^{\mathrm{th}}$--$95^{\mathrm{th}}$-percentile band across
all 1000 iterations, with the vancomycin and linezolid dosing windows
overlaid.
\item \textbf{\texttt{mc\_ec50\_lzd\_sweep\_outcomes.png}} -- four scatter
panels (peak $S_b$, peak $R_b$, peak $S_{res}$, final $R_b$) plotted against
$\log_{10}(EC_{50,L})$ for each of the 1000 samples, with the baseline
$EC_{50,L}$ marked.
\item \textbf{\texttt{mc\_ec50\_lzd\_switch.png}} -- a two-panel figure: the
top panel is the continuous deterministic switch curve (final $R_b$/$R_{res}$
vs.\ $EC_{50,L}$) with \texttt{ESCAPE\_THRESH} marked and the suppressed/
escape regimes shaded; the bottom panel is a histogram of the 1000 sampled
$EC_{50,L}$ values, split and colored by each patient's actual simulated
outcome (suppressed vs.\ escape) rather than by the deterministic threshold
alone.
\end{enumerate}

\end{document}
